\subsection{Latar Belakang}
Kebutuhan akses internet berkecepatan tinggi menjadi pilar fundamental dalam mendukung aktivitas tridharma perguruan tinggi di Universitas Logistik dan Bisnis Internasional (ULBI). Pengelolaan akses Wi-Fi kampus umumnya memanfaatkan mekanisme \textit{captive portal}, di mana lalu lintas data pengguna diarahkan secara paksa ke halaman web autentikasi sebelum izin penjelajahan internet diberikan.

Meskipun efektif untuk manajemen hak akses dan kuota, penerapan \textit{captive portal} menimbulkan friksi teknis yang signifikan bagi mahasiswa dan dosen:
\begin{enumerate}
    \item \textbf{Pengisian Berulang:} Mahasiswa berpindah ruangan perkuliahan atau laboratorium sering kali mengalami pemutusan sesi (\textit{session timeout}), sehingga wajib membuka peramban web dan memasukkan kembali nama pengguna serta kata sandi secara manual.
    \item \textbf{Kegagalan Redirection:} Pada sejumlah sistem operasi Android modern, fitur \textit{captive portal detection} terkadang tidak membuka antarmuka peramban secara otomatis atau memunculkan peringatan ketiadaan koneksi internet tanpa opsi langsung masuk portal.
    \item \textbf{Keamanan Kredensial:} Pengetikan kredensial di tempat umum rentan terhadap intipan visual (\textit{shoulder surfing}) maupun kekeliruan penyimpanan kata sandi pada peramban publik.
\end{enumerate}

Untuk mengatasi tantangan tersebut, Proyek 2 ini mengembangkan aplikasi perangkat bergerak \textbf{Coneckin} berbasis Android \textit{native} dengan nama paket \texttt{coneckin.id.biz.id}. Coneckin dibangun untuk memberikan kemudahan koneksi satu ketukan (\textit{one-tap connection}) dengan tingkat keamanan tinggi melalui enkripsi lokal berbasis perangkat keras (\textit{hardware-backed cryptography}).

\subsection{Profil Kompetensi Capstone: Software Engineer}
Sesuai dengan ketentuan kurikulum Proyek Capstone Sarjana Terapan Teknik Informatika ULBI dan panduan proyek pada \url{https://ll.my.id/instruction/}, Proyek 2 berfokus pada pembuktian profil kompetensi \textbf{Software Engineer}. Pada tahap ini, mahasiswa diwajibkan untuk:
\begin{enumerate}
    \item Membangun fitur perangkat lunak secara utuh dan \textit{end-to-end} (analisis kebutuhan $\rightarrow$ perancangan arsitektur $\rightarrow$ implementasi $\rightarrow$ integrasi).
    \item Menulis kode pengujian mandiri (\textit{unit testing} dan \textit{integration testing}) untuk memverifikasi kebenaran logika bisnis.
    \item Mengintegrasikan \textit{Continuous Integration / Continuous Deployment} (CI/CD) pada GitHub Actions untuk menjalankan build dan pengujian secara otomatis setiap kali ada pembaruan kode.
    \item Menghasilkan artefak rilis berupa berkas biner (APK/AAB) serta dokumentasi terdistribusi.
\end{enumerate}

\subsection{Rumusan Masalah}
Berdasarkan latar belakang yang diuraikan, rumusan masalah dalam Proyek 2 ini adalah:
\begin{enumerate}
    \item Bagaimana merancang arsitektur aplikasi Android yang mampu mendeteksi dan mengotomasi proses login \textit{captive portal} Wi-Fi kampus ULBI secara andal dan cepat?
    \item Bagaimana mengamankan penyimpanan kredensial akun multipengguna pada penyimpanan lokal perangkat Android tanpa mengirimkan kata sandi ke server pihak ketiga?
    \item Bagaimana mengimplementasikan integrasi antarmuka \textit{Quick Settings Tile} pada panel notifikasi Android untuk memfasilitasi akses login satu ketukan?
    \item Bagaimana menyusun rangkaian pengujian otomatis dan pipeline CI/CD untuk memastikan kualitas serta kestabilan rilis aplikasi?
\end{enumerate}

\subsection{Batasan Masalah}
Batasan masalah dalam pelaksanaan proyek ini meliputi:
\begin{enumerate}
    \item Aplikasi dikembangkan untuk sistem operasi Android dengan target SDK versi 34 (Android 14) dan batas minimum SDK versi 29 (Android 10).
    \item Pengujian otomasi login difokuskan pada infrastruktur \textit{captive portal} Wi-Fi kampus ULBI dengan \textit{hostname} \texttt{ypbpi.hotspot}.
    \item Kredensial pengguna disimpan secara lokal terenkripsi di dalam perangkat dan tidak pernah disinkronkan ke server eksternal demi privasi pengguna.
\end{enumerate}

\subsection{Tujuan Proyek}
Tujuan yang dicapai melalui proyek ini adalah:
\begin{enumerate}
    \item Menghasilkan aplikasi Android Coneckin yang berfungsi secara penuh untuk melakukan otomasi login \textit{captive portal} kampus.
    \item Menerapkan arsitektur keamanan dengan modul \textit{CryptoManager} yang memanfaatkan Android Keystore dan autentikasi biometrik.
    \item Mengimplementasikan komponen \textit{Quick Settings Tile} pada antarmuka Android.
    \item Membangun alur CI/CD pada GitHub Actions yang mencakup build otomatis, pengujian unit, serta rilis otomatis APK dan dokumen PDF.
\end{enumerate}

\subsection{Manfaat Proyek}
Manfaat dari penelitian dan pengembangan ini adalah:
\begin{enumerate}
    \item \textbf{Bagi Mahasiswa/Pengguna:} Menghemat waktu akses jaringan Wi-Fi di lingkungan kampus dan meningkatkan kenyamanan aktivitas akademik.
    \item \textbf{Bagi Institusi (ULBI):} Menyediakan portofolio produk perangkat lunak terapan yang siap pakai karya mahasiswa program studi Sarjana Terapan Teknik Informatika.
    \item \textbf{Bagi Pengembang:} Memenuhi capaian pembelajaran mata kuliah Proyek 2 serta mengasah kompetensi rekayasa perangkat lunak modern.
\end{enumerate}
